\section{The randomized lower bound}
\label{sec:classical}

We compare a distribution with one zero in every column to a coupled
distribution in which one uniformly chosen column is made a winner.
An algorithm can distinguish them only by finding the zero that was
removed or by reaching the newly filled certificate cell.  The first
event requires searching many columns; the second requires either many
Forrelation queries in the chosen column or a successful guess of its
entire address.

\subsection{Independent answers under adaptive queries}

We first record the precise amplification fact needed for address
guessing.  We use the capped direct-product argument developed by
Drucker~\cite[Section~1.3 and Lemmas~3.1--3.2]{Drucker12}, following
earlier work of Shaltiel~\cite{Shaltiel03}.  Its query limit is imposed
separately on each input, while queries to different inputs may be
interleaved adaptively.

\Needspace{12\baselineskip}
\begin{lemma}\label{lem:adaptive-product}
Let $P\colon\mathcal D\to\bits$, where $\mathcal D\subseteq\bits^n$, let
$\mu$ be a distribution on $\mathcal D$, and let $L$ be a nonnegative
integer.  Suppose every deterministic algorithm making at most $L$
queries computes $P$ under $\mu$ with success probability at most $\beta$.
If $U_1,\ldots,U_r$ are independent samples from $\mu$, every classical
algorithm making at most $L$ queries to each $U_h$ has probability at most
$\beta^r$ of outputting
\[
 (P(U_1),\ldots,P(U_r)).
\]
The algorithm may be randomized, choose its queries adaptively among the
$r$ inputs, and receive arbitrary additional information independent of
$(U_1,\ldots,U_r)$.
\end{lemma}

\begin{proof}
Fix the additional information and all internal random coins, leaving a
deterministic algorithm.  At any transcript $\tau$ of positive
probability, the conditional law of $(U_1,\ldots,U_r)$ is a product of
their separate conditional laws.  Indeed, the event of obtaining $\tau$
requires exactly the observed bit values in each input.  Any tuple
satisfying these requirements follows the same query schedule, since
the next query is determined by the preceding answers.  Thus this event
is a Cartesian product of constraints on the individual inputs.

For each $h$, let $q_h$ be the number of queries already made to $U_h$,
and let $V_h(\tau)$ be the largest success probability for computing
$P(U_h)$ from its current conditional law using at most $L-q_h$ further
queries.  This maximum is attained because the domain and the remaining
query budget are finite.  If the next query addresses $U_h$, optimality
gives
\begin{equation}\label{eq:bellman-query}
 \mathbb E[V_h(\tau')\mid\tau]\leq V_h(\tau),
\end{equation}
where $\tau'$ includes its answer: making that particular query and
then using an optimal continuation is one allowed strategy in the
definition of $V_h(\tau)$.  The conditional laws and budgets of the
other inputs do not change.  Hence
\[
 \mathbb E\!\left[
   \prod_{j=1}^r V_j(\tau')\,\middle|\,\tau
 \right]
 \leq \prod_{j=1}^r V_j(\tau).
\]
The product is therefore a supermartingale over the successive queries.
There are at most $rL$ such queries; a run that stops earlier can be
extended by steps leaving all values unchanged.

At the final transcript, the conditional probability that the algorithm's
chosen answer vector is correct is at most
\[
 \prod_{h=1}^r
   \max_{b\in\bits}\Pr(P(U_h)=b\mid\tau)
 \leq \prod_{h=1}^r V_h(\tau).
\]
The first inequality uses the product conditional law, and the second
uses the option of making no further queries.  Initially every $V_h$
is at most $\beta$.  Taking expectations proves the bound $\beta^r$.
Averaging over the fixed information and coins proves the full statement.
\end{proof}

Return to $P=P_{k,m}$ and $L$ from \eqref{eq:L-def}, with $m$ large
enough that $1\leq L<\R(P)$.  There are finitely many deterministic
decision trees of depth at most $L$ on this finite domain.  A randomized
algorithm of that query depth is a distribution on these trees. Finite
minimax, in its query-complexity formulation~\cite{Yao77}, applied to the
payoff that the answer is correct, gives
\[
 \begin{aligned}
 &\max_{\text{distributions on depth-$L$ trees}}
       \min_{x\in\operatorname{dom}P}\Pr(\text{correct on }x)\\
 &\hspace{2em}=
 \min_{\mu}\max_{\text{depth-$L$ trees }T}
       \Pr_{x\sim\mu}(T(x)=P(x)).
 \end{aligned}
\]
The left-hand side is less than $2/3$, since its maximum is attained
and $L<\R(P)$.  Consequently there is a distribution $\mu$ on
$\operatorname{dom}P$ for which every such deterministic tree has
success probability at most $2/3$.  Lemma~\ref{lem:adaptive-product}
therefore bounds the success probability of guessing $r$ independent
answers by $(2/3)^r$.

\subsection{Coupled zero and one inputs}

\begin{proposition}\label{prop:classical}
For each fixed integer $k\geq2$ and every sufficiently large power of
four $m$, the total function of Section~\ref{sec:construction} satisfies
\begin{equation}\label{eq:classical-lower}
 \R(F_{k,m})>
 G:=\left\lfloor\frac{mL}{100}\right\rfloor,
 \qquad
 L=\left\lfloor\frac{m^{1-1/k}}{(\log_2m)^2}\right\rfloor.
\end{equation}
\end{proposition}

\begin{proof}
Generate a \emph{dummy input} $Y^{(0)}$ as follows.  Independently for
each $i\in[m]$, choose $J_i$ uniformly from $[m]$ and put a single zero
at position $J_i$ in $X_i$, setting its other bits to 1.  Draw all
$A_{i,h}$ independently from $\mu$, independently also of the $J_i$.
Set all cells $B_{i,s}$ to one fixed blank bit pattern.  Since each $X_i$
has a zero, $F_{k,m}(Y^{(0)})=0$.

Choose $I$ uniformly from $[m]$, independently of the entire dummy
input.  Form a \emph{planted input} $Y^{(1)}$ by setting $X_I=1^m$ and
filling its true-address cell $B_{I,s_I}$ with the correct circuit
transcripts and the pointers $p_{I,s_I,j}=J_j$ for $j\ne I$.
All $A$-strings are in the domain of $P$, so $s_I$ is well defined.
Column $I$ is a winner, and therefore $F_{k,m}(Y^{(1)})=1$.

Fix a deterministic algorithm making at most $G$ queries, and consider
its run on $Y^{(0)}$.  The coupled inputs can differ only at the bit
$X_I[J_I]$ and inside the cell $B_{I,s_I}$.  Until one of these locations
is queried, the two executions have identical answers and hence
identical subsequent queries.  It follows that the difference between
their output-1 probabilities is at most
\begin{equation}\label{eq:coupling-events}
 \Pr(\text{the dummy run queries }X_I[J_I])
 +\Pr(\text{the dummy run queries any bit of }B_{I,s_I}).
\end{equation}
Counting every visit to the indicated cell is an upper bound even if
the queried bit happens to agree with its blank value.

\paragraph{Finding the removed zero.}
Let $Z$ be the number of columns whose zero is discovered in the dummy
run.  Before the zero in a column is found, its location is uniform
among the positions in that column not yet queried.  This remains true
conditional on the full adaptive transcript: the other $X$-strings and
all $A$-strings are independent of its zero location, and the blank
cells supply no information about that location.

Call a discovery early if fewer than $m/2$ distinct positions in its
column were queried before it.  Conditional on the preceding transcript,
each query has probability at most $2/m$ of an early discovery.  Thus
the expected number of early discoveries is at most $2G/m$.  Each
remaining discovery occurs in a column that receives at least $m/2$
queries, and there are at most $2G/m$ such columns.  Hence
\[
 \mathbb E Z\leq\frac{4G}{m}.
\]
The dummy run is independent of $I$, so
\begin{equation}\label{eq:zero-hit}
 \Pr(\text{the dummy run queries }X_I[J_I])
 =\frac{\mathbb E Z}{m}\leq\frac{4G}{m^2}.
\end{equation}

\paragraph{Reaching the planted cell.}
Call a column heavy if, over the entire dummy run, the algorithm makes
more than $L$ queries in total to its strings $A_{i,1},\ldots,A_{i,r}$.
There are at most $G/L$ heavy columns.  Independence of $I$ gives
\begin{equation}\label{eq:heavy-probability}
 \Pr(I\text{ is heavy})\leq\frac{G}{mL}.
\end{equation}
For the other columns we use Lemma~\ref{lem:adaptive-product}.  The
following simulation bounds a joint event and does not condition the
input distribution on a column being nonheavy.

Fix $i\in[m]$ and a time $t\in\{1,\ldots,G\}$.  Give a simulator all
the dummy input outside $A_{i,1},\ldots,A_{i,r}$ for free.  This
information is independent of those $r$ strings.  The simulator runs
the deterministic algorithm up to its proposed $t$th query, but aborts
before making an $(L+1)$st query to these $r$ strings in total.  If it
reaches time $t$ and the proposed query is to a cell $B_{i,s}$, it outputs
the address $s$ as its guess for the vector of $r$ answers.  Otherwise,
including if the simulated run stops or aborts, it outputs any fixed
vector.

The simulator makes at most $L$ queries in total to the target strings,
and in particular at most $L$ to each one.  By
Lemma~\ref{lem:adaptive-product}, its probability of guessing $s_i$
correctly is at most $(2/3)^r$.  Whenever the original dummy run queries
$B_{i,s_i}$ at time $t$ and column $i$ is nonheavy, this simulator has
not aborted and outputs the correct address.  Thus
\begin{equation}\label{eq:capped-address}
 \Pr\bigl(i\text{ is nonheavy and the $t$th query is in }B_{i,s_i}\bigr)
 \leq (2/3)^r.
\end{equation}
Union bounding over $t$ and averaging over the independent uniform
choice of $I$, then adding \eqref{eq:heavy-probability}, yields
\begin{equation}\label{eq:cell-hit}
 \Pr(\text{the dummy run queries }B_{I,s_I})
 \leq\frac{G}{mL}+G(2/3)^r.
\end{equation}

Equations~\eqref{eq:coupling-events}, \eqref{eq:zero-hit}, and
\eqref{eq:cell-hit} bound the difference of output probabilities by
\begin{equation}\label{eq:classical-distinguishing}
 \frac{4G}{m^2}+\frac{G}{mL}+G(2/3)^r.
\end{equation}
For all sufficiently large $m$, we have $1\leq L\leq m$ and
$G\leq mL/100$.  The first two terms are therefore at most $4/100$
and $1/100$, respectively.  Moreover, \eqref{eq:r-def} gives
\[
 G(2/3)^r
 \leq\frac1{100}m^{\,2-20\log_2(3/2)}=o(1).
\]
In particular, \eqref{eq:classical-distinguishing} is strictly less
than $1/3$ for large $m$.

The same bound holds for a randomized algorithm of worst-case query
cost $G$, by fixing its coins and then averaging.  Such an algorithm,
if it had worst-case error at most $1/3$, would output 1 with probability
at most $1/3$ on $Y^{(0)}$ and at least $2/3$ on $Y^{(1)}$.  Their
output-1 probabilities would differ by at least $1/3$, contradicting
the preceding estimate.  Hence $\R(F_{k,m})>G$.
\end{proof}

For fixed $k$, Proposition~\ref{prop:classical} implies
\[
 \R(F_{k,m})=
 \Omega_k\!\left(\frac{m^{2-1/k}}{(\log m)^2}\right).
\]
The heavy-column estimate pays for learning an address, while the
independent-answer estimate makes guessing an unlearned address
negligible even over all $G$ adaptive queries.
